\documentclass[12pt, a4paper]{article}

% ============================================================
%  Packages
% ============================================================
\usepackage[T1]{fontenc}
\usepackage[utf8]{inputenc}
\usepackage[french]{babel}
\usepackage[
  top=2.5cm, bottom=2.5cm,
  left=3cm,  right=3cm
]{geometry}
\usepackage{amsmath, amssymb, amsthm}
\usepackage{booktabs}
\usepackage{hyperref}
\hypersetup{
  colorlinks = true,
  linkcolor  = blue!70!black,
  urlcolor   = blue!70!black,
  citecolor  = blue!70!black,
  pdftitle   = {Fondements Mathématiques des Graph Neural Networks Spectraux},
  pdfauthor  = {S. Oussama},
}

% ============================================================
%  Environnements théorème
% ============================================================
\newtheorem{theorem}{Théorème}[section]
\newtheorem{lemma}[theorem]{Lemme}
\newtheorem{corollary}[theorem]{Corollaire}
\newtheorem{proposition}[theorem]{Proposition}

\theoremstyle{definition}
\newtheorem{definition}[theorem]{Définition}
\newtheorem{example}[theorem]{Exemple}
\newtheorem{remark}[theorem]{Remarque}

% ============================================================
%  Macros
% ============================================================
\newcommand{\R}{\mathbb{R}}
\newcommand{\N}{\mathbb{N}}
\newcommand{\calG}{\mathcal{G}}
\newcommand{\calV}{\mathcal{V}}
\newcommand{\calE}{\mathcal{E}}
\newcommand{\mat}[1]{\mathbf{#1}}
\newcommand{\A}{\mat{A}}
\newcommand{\D}{\mat{D}}
\newcommand{\I}{\mat{I}}
\newcommand{\L}{\mat{L}}
\newcommand{\Lt}{\tilde{\mat{L}}}
\newcommand{\Lam}{\mat{\Lambda}}
\newcommand{\U}{\mat{U}}
\newcommand{\x}{\mathbf{x}}
\newcommand{\y}{\mathbf{y}}
\newcommand{\norm}[1]{\left\lVert #1 \right\rVert}

% ============================================================
%  Page de titre
% ============================================================
\title{%
  \vspace{-1cm}
  {\LARGE\bfseries Fondements Mathématiques des Graph Neural Networks Spectraux}\\[0.8em]
  {\large De la Théorie Spectrale des Graphes à l'Approximation de Tchebychev}\\[0.5em]
  {\normalsize Projet \texttt{spectral-gnn} --- \texttt{governor08/spectral-gnn}}
}
\author{S. Oussama \\ \small Software Engineer \& Math Enthusiast}
\date{Avril 2026}

% ============================================================
\begin{document}
% ============================================================

\maketitle
\thispagestyle{empty}

\begin{abstract}
Ce document établit les fondements mathématiques rigoureux du projet \texttt{spectral-gnn},
une implémentation \emph{from scratch} (sans PyTorch Geometric) d'un réseau de neurones
convolutif sur graphe basé sur les polynômes de Tchebychev (ChebNet,
Defferrard \emph{et al.}, 2016).

La progression suit une logique de \emph{problématique $\to$ démonstration} :
nous partons du problème de la définition d'une convolution sur un domaine non-euclidien
(Section~\ref{sec:laplacian}), nous montrons comment la diagonalisation du Laplacien
fournit une réponse exacte mais cubiquement coûteuse (Section~\ref{sec:spectral}),
puis nous démontrons comment l'approximation polynomiale de Tchebychev ramène
cette complexité à un régime linéaire en le nombre d'arêtes (Section~\ref{sec:chebnet}).

L'implémentation résultante, construite sur des tenseurs PyTorch bruts,
atteint \textbf{78.2\% de précision} sur le dataset Cora en classification
semi-supervisée de nœuds --- résultat conforme à celui rapporté dans
Defferrard \emph{et al.} (2016), validant ainsi la correction de chaque étape
mathématique décrite dans ce document.
\end{abstract}

\newpage
\tableofcontents
\newpage

% ============================================================
\section{Le Laplacien de Graphe}
\label{sec:laplacian}
% ============================================================

\subsection{Problématique : comment définir une structure différentielle sur un graphe ?}

En traitement du signal classique, la convolution entre un signal $f$ et un filtre $g$
est définie via la transformation de Fourier, laquelle repose sur les fonctions propres
de l'opérateur laplacien $\Delta$ sur $\R^n$.
Sur un graphe fini, il n'existe \emph{a priori} aucune notion de dérivée directionnelle.
L'idée fondatrice de la théorie spectrale des graphes est de construire un analogue discret
du laplacien continu, dont les vecteurs propres joueront le rôle des fréquences de Fourier.

\subsection{Graphes et matrices associées}

\begin{definition}[Graphe non orienté pondéré]
Un graphe $\calG = (\calV, \calE, w)$ est la donnée :
\begin{itemize}
  \item d'un ensemble de $N$ sommets (noeuds) $\calV = \{1, \ldots, N\}$,
  \item d'un ensemble d'arêtes $\calE \subseteq \calV \times \calV$,
  \item d'une fonction de poids $w : \calE \to \R_{>0}$.
\end{itemize}
Dans la suite, on considère uniquement des graphes \textbf{non orientés}
($\{i,j\} \in \calE \Rightarrow \{j,i\} \in \calE$) et \textbf{sans boucle}
($\{i,i\} \notin \calE$).
\end{definition}

\begin{definition}[Matrice d'adjacence]
La \emph{matrice d'adjacence} $\A \in \R^{N \times N}$ est définie par :
\[
  A_{ij} =
  \begin{cases}
    w(i,j) & \text{si } \{i,j\} \in \calE, \\
    0       & \text{sinon.}
  \end{cases}
\]
Par la contrainte de non-orientation, $\A$ est \textbf{symétrique} : $\A = \A^\top$.
\end{definition}

\begin{definition}[Matrice de degrés]
La \emph{matrice de degrés} $\D \in \R^{N \times N}$ est la matrice diagonale :
\[
  D_{ii} = d_i := \sum_{j=1}^{N} A_{ij}, \qquad D_{ij} = 0 \text{ pour } i \neq j.
\]
$d_i$ est appelé le \emph{degré} du nœud $i$.
\end{definition}

\subsection{Construction du Laplacien normalisé}

\begin{definition}[Laplacien combinatoire]
Le \emph{Laplacien combinatoire} (ou non normalisé) du graphe $\calG$ est :
\[
  \L_{\mathrm{comb}} = \D - \A.
\]
\end{definition}

\begin{remark}
Le Laplacien combinatoire généralise l'opérateur discret de diffusion :
pour un signal $\x \in \R^N$, on a
$(\L_{\mathrm{comb}} \x)_i = d_i x_i - \sum_{j \sim i} w(i,j)\, x_j$,
ce qui mesure la différence entre la valeur de $\x$ en $i$ et la somme pondérée
de ses voisins --- analogue discret de $-\Delta f$.
\end{remark}

\begin{definition}[Laplacien normalisé symétrique]
\label{def:normalized_laplacian}
Le \emph{Laplacien normalisé symétrique} est :
\[
  \boxed{\L = \I - \D^{-1/2} \A \D^{-1/2}}
\]
où $\D^{-1/2}$ est la matrice diagonale dont les entrées valent
$D^{-1/2}_{ii} = d_i^{-1/2}$ pour tout nœud $i$ de degré non nul.
\end{definition}

\begin{remark}
On peut vérifier que $\L_{ij}$ prend les valeurs suivantes :
\[
  L_{ij} =
  \begin{cases}
    1 & \text{si } i = j \text{ et } d_i \neq 0, \\
    -\dfrac{w(i,j)}{\sqrt{d_i d_j}} & \text{si } \{i,j\} \in \calE, \\
    0 & \text{sinon.}
  \end{cases}
\]
\end{remark}

\subsection{Propriétés spectrales du Laplacien normalisé}

\begin{theorem}[Spectre du Laplacien normalisé]
\label{thm:spectrum}
Soit $\calG$ un graphe connexe à $N$ nœuds sans boucle.
La matrice $\L = \I - \D^{-1/2} \A \D^{-1/2}$ est :
\begin{enumerate}
  \item \textbf{Symétrique réelle}, donc diagonalisable dans une base orthonormée.
  \item \textbf{Semidéfinie positive} : toutes ses valeurs propres vérifient $\lambda \geq 0$.
  \item \textbf{Bornée supérieurement} : toutes ses valeurs propres vérifient $\lambda \leq 2$.
\end{enumerate}
En résumé : $\sigma(\L) \subset [0, 2]$.
\end{theorem}

\begin{proof}
\textbf{(1) Symétrie.}
Puisque $\A$ est symétrique et $\D^{-1/2}$ est diagonale réelle,
on a $(\D^{-1/2} \A \D^{-1/2})^\top = \D^{-1/2} \A^\top \D^{-1/2} = \D^{-1/2} \A \D^{-1/2}$.
Donc $\L^\top = \I^\top - (\D^{-1/2} \A \D^{-1/2})^\top = \L$.

\textbf{(2) Semidéfinition positive.}
Pour tout $\x \in \R^N$, on calcule directement la forme quadratique associée à $\L$ :
\[
  \x^\top \L \x
  = \sum_{\{i,j\} \in \calE} w(i,j)
    \left( \frac{x_i}{\sqrt{d_i}} - \frac{x_j}{\sqrt{d_j}} \right)^2 \geq 0.
\]
Cette forme quadratique est une somme de carrés, donc $\L \succeq 0$.
En particulier, toutes les valeurs propres satisfont $\lambda \geq 0$.

\textbf{(3) Borne supérieure $\lambda \leq 2$ et borne inférieure $\lambda \geq 0$.}
Notons $\hat{\A} = \D^{-1/2} \A \D^{-1/2}$, de sorte que $\L = \I - \hat{\A}$.
Si $\mu$ est une valeur propre de $\hat{\A}$, alors $\lambda = 1 - \mu$ est valeur propre
de $\L$.

\emph{Borne supérieure ($\mu \leq 1$, i.e. $\lambda \geq 0$).}
Ceci découle directement de la semidéfinition positive de $\L$ démontrée au point (2) :
$\lambda = 1 - \mu \geq 0$ implique $\mu \leq 1$.

\emph{Borne inférieure ($\mu \geq -1$, i.e. $\lambda \leq 2$).}
Soit $v \in \R^N$ un vecteur unitaire ($\norm{v} = 1$). On minore la forme quadratique
de $\hat{\A}$ en utilisant l'inégalité arithmético-géométrique ($AM$-$GM$) :
pour tous réels $a, b$, on a $ab \geq -\frac{1}{2}(a^2 + b^2)$.
Appliquée à $a = v_i / \sqrt{d_i}$ et $b = v_j / \sqrt{d_j}$ :
\[
  \frac{v_i v_j}{\sqrt{d_i d_j}} \geq -\frac{1}{2}\left(\frac{v_i^2}{d_i} + \frac{v_j^2}{d_j}\right).
\]
En sommant sur toutes les arêtes $\{i,j\} \in \calE$ (avec $w(i,j) = 1$ pour simplifier) :
\[
  v^\top \hat{\A}\, v
  = \sum_{\{i,j\} \in \calE} \frac{v_i v_j}{\sqrt{d_i d_j}}
  \geq -\frac{1}{2} \sum_{\{i,j\} \in \calE} \left(\frac{v_i^2}{d_i} + \frac{v_j^2}{d_j}\right).
\]
Or, en regroupant par nœud $i$ :
\[
  \sum_{\{i,j\} \in \calE} \frac{v_i^2}{d_i} = \sum_{i=1}^N v_i^2 \cdot \frac{d_i}{d_i} = \sum_{i=1}^N v_i^2 = \norm{v}^2 = 1,
\]
et de même pour le terme en $v_j^2 / d_j$. Donc :
\[
  v^\top \hat{\A}\, v \geq -\frac{1}{2}(1 + 1) = -1.
\]
Ainsi $\mu = v^\top \hat{\A}\, v \geq -1$ pour tout vecteur propre unitaire $v$,
ce qui donne $\lambda = 1 - \mu \leq 2$.

En conclusion, $\mu(\hat{\A}) \in [-1, 1]$ et $\lambda(\L) = 1 - \mu(\hat{\A}) \in [0, 2]$.
\end{proof}

\begin{corollary}
\label{cor:eigenvector_zero}
Le spectre de $\L$ est de la forme $0 = \lambda_1 \leq \lambda_2 \leq \cdots \leq \lambda_N \leq 2$.
La valeur propre $\lambda_1 = 0$ correspond au vecteur propre
$u_1 = \D^{1/2} \mathbf{1} / \norm{\D^{1/2} \mathbf{1}}$.
Le graphe est connexe si et seulement si $\lambda_2 > 0$.
\end{corollary}

\begin{proof}
Vérifions directement que $\L \cdot \D^{1/2}\mathbf{1} = \mathbf{0}$ :
\begin{align*}
  \D^{-1/2} \A \D^{-1/2} \cdot \D^{1/2}\mathbf{1}
  &= \D^{-1/2} \A \mathbf{1}
   = \D^{-1/2} (d_1, \ldots, d_N)^\top \\
  &= \D^{-1/2} \D \mathbf{1}
   = \D^{1/2} \mathbf{1}.
\end{align*}
Donc $(\I - \D^{-1/2} \A \D^{-1/2})\, \D^{1/2}\mathbf{1}
      = \D^{1/2}\mathbf{1} - \D^{1/2}\mathbf{1} = \mathbf{0}$.
Le vecteur propre associé à $\lambda_1 = 0$ est bien $u_1 = \D^{1/2}\mathbf{1}/\norm{\D^{1/2}\mathbf{1}}$.
\end{proof}

% ============================================================
\section{La Convolution Spectrale Stricte}
\label{sec:spectral}
% ============================================================

\subsection{Problématique : comment filtrer un signal sur graphe ?}

En traitement du signal \emph{discret}, la convolution circulaire de deux signaux
$x, h \in \R^N$ est diagonale dans la base de la transformée de Fourier discrète (DFT).
Plus précisément, les vecteurs de base
\[
  e_k = \frac{1}{\sqrt{N}}\bigl(1,\, \omega^k,\, \omega^{2k},\, \ldots,\, \omega^{(N-1)k}\bigr)^\top,
  \qquad \omega = e^{2\pi i / N},
\]
sont exactement les vecteurs propres de la matrice de décalage circulaire.
La convolution revient alors à une multiplication point par point ({\em pointwise})
dans ce domaine fréquentiel : $\widehat{x * h} = \hat{x} \cdot \hat{h}$.

Sur un graphe quelconque, la matrice de décalage circulaire est remplacée par le
Laplacien normalisé $\L$, dont les vecteurs propres $\U = [u_1 | \cdots | u_N]$
jouent le rôle des vecteurs de Fourier.
On définit ainsi une transformée de Fourier sur graphe, et la convolution spectrale
généralise l'opération classique à des domaines irréguliers.

\subsection{Décomposition spectrale du Laplacien}

Puisque $\L$ est symétrique réelle (Théorème~\ref{thm:spectrum}),
le théorème spectral réel garantit sa diagonalisation dans une base orthonormée :

\begin{theorem}[Décomposition spectrale]
\label{thm:diag}
Il existe une matrice orthogonale $\U \in \R^{N \times N}$ ($\U^\top \U = \U \U^\top = \I$)
et une matrice diagonale $\Lam = \mathrm{diag}(\lambda_1, \ldots, \lambda_N)$
avec $0 \leq \lambda_1 \leq \cdots \leq \lambda_N \leq 2$, telles que :
\[
  \boxed{\L = \U \Lam \U^\top}
\]
Les colonnes $u_1, \ldots, u_N$ de $\U$ sont les \emph{vecteurs propres} de $\L$,
qui constituent la \textbf{base de Fourier sur graphe}.
\end{theorem}

\subsection{Transformée de Fourier sur graphe}

\begin{definition}[Transformée de Fourier sur graphe]
Pour un signal $\x \in \R^N$ défini sur les nœuds du graphe, la
\emph{transformée de Fourier sur graphe} et son inverse sont :
\[
  \hat{\x} = \U^\top \x \in \R^N
  \qquad \text{et} \qquad
  \x = \U \hat{\x}.
\]
$\hat{x}_k$ est l'amplitude de la composante fréquentielle $u_k$ (associée à $\lambda_k$).
\end{definition}

\subsection{Convolution spectrale et filtrage}

\begin{definition}[Convolution spectrale sur graphe]
\label{def:spectral_conv}
La convolution d'un signal $\x$ avec un filtre spectral $g_\theta$ est définie par :
\[
  \boxed{g_\theta * \x = \U\, g_\theta(\Lam)\, \U^\top \x}
\]
où $g_\theta(\Lam) = \mathrm{diag}\bigl(g_\theta(\lambda_1), \ldots, g_\theta(\lambda_N)\bigr)$
est le filtre appliqué diagonalement dans le domaine spectral,
et $\theta \in \R^N$ est le vecteur des paramètres apprenables.
\end{definition}

\begin{remark}[Interprétation]
L'opération se décompose en trois étapes :
\begin{enumerate}
  \item \textbf{Analyse} : $\hat{\x} = \U^\top \x$ --- projection sur la base de Fourier.
  \item \textbf{Filtrage} : $\hat{\x}' = g_\theta(\Lam)\, \hat{\x}$ --- multiplication pointwise des coefficients fréquentiels.
  \item \textbf{Synthèse} : $\y = \U \hat{\x}'$ --- reconstruction dans le domaine nodal.
\end{enumerate}
\end{remark}

\subsection{Complexité de l'approche exacte}

\begin{proposition}[Complexité de la convolution spectrale stricte]
\label{prop:complexity_strict}
L'évaluation de $g_\theta * \x = \U g_\theta(\Lam) \U^\top \x$ a une complexité :
\[
  \mathcal{O}(N^3)
\]
en temps (pour la diagonalisation de $\L$) et $\mathcal{O}(N^2)$ en espace
(pour le stockage de $\U$).
\end{proposition}

\begin{proof}
La décomposition en valeurs propres d'une matrice $N \times N$ par l'algorithme QR
(ou LAPACK \texttt{dsyev}) a une complexité $\mathcal{O}(N^3)$.
Une fois $\U$ calculé, le produit $\U^\top \x$ coûte $\mathcal{O}(N^2)$,
le produit $g_\theta(\Lam) \hat{\x}$ coûte $\mathcal{O}(N)$
(multiplication diagonale), et $\U \hat{\x}'$ coûte $\mathcal{O}(N^2)$.
Le goulot d'étranglement est donc la diagonalisation : $\mathcal{O}(N^3)$.

Pour des graphes réels (Cora : $N = 2708$, ogbn-arxiv : $N = 169\,343$),
cette complexité est prohibitive.
\end{proof}

\begin{remark}[Limitation supplémentaire]
Outre la complexité, le filtre $g_\theta \in \R^N$ possède autant de paramètres
que de nœuds dans le graphe, ce qui le rend non transférable d'un graphe à l'autre
et extrêmement susceptible au surapprentissage.
\end{remark}

% ============================================================
\section{L'Approximation de Tchebychev (ChebNet)}
\label{sec:chebnet}
% ============================================================

\subsection{Problématique : peut-on éviter la diagonalisation ?}

L'approche spectrale stricte (Section~\ref{sec:spectral}) souffre de deux défauts :
\begin{itemize}
  \item Complexité $\mathcal{O}(N^3)$ pour la diagonalisation.
  \item Filtre non localisé : $g_\theta(\Lam)$ peut mélanger des nœuds arbitrairement distants.
\end{itemize}
L'idée de Hammond, Vandergheynst \& Gribonval (2011), reprise par Defferrard,
Bresson \& Vandergheynst (2016) dans ChebNet, est d'approcher $g_\theta(\Lam)$
par un polynôme en $\Lam$ d'ordre $K$, ce qui permet de \emph{remplacer}
la diagonalisation par des multiplications matricielles creuses.

\subsection{Rescaling du spectre}

\begin{definition}[Laplacien normalisé rescalé]
\label{def:rescaled_laplacian}
Afin que les polynômes de Tchebychev soient évalués sur leur domaine d'orthogonalité
$[-1, 1]$, on rescale $\L$ en :
\[
  \boxed{\Lt = \frac{2}{\lambda_{\max}} \L - \I}
\]
où $\lambda_{\max} = \max_k \lambda_k$ est la plus grande valeur propre de $\L$.
\end{definition}

\begin{proposition}
Les valeurs propres de $\Lt$ appartiennent à $[-1, 1]$.
\end{proposition}

\begin{proof}
Si $\lambda_k \in [0, \lambda_{\max}]$, alors
$\tilde{\lambda}_k = \frac{2\lambda_k}{\lambda_{\max}} - 1 \in [-1, 1]$.
\end{proof}

\subsection{Polynômes de Tchebychev}

\begin{definition}[Polynômes de Tchebychev de première espèce]
Les \emph{polynômes de Tchebychev} $(T_k)_{k \geq 0}$ sont définis sur $[-1, 1]$
par la relation de récurrence à trois termes :
\[
  \boxed{
    T_0(x) = 1, \quad
    T_1(x) = x, \quad
    T_k(x) = 2x\, T_{k-1}(x) - T_{k-2}(x), \quad k \geq 2.
  }
\]
\end{definition}

\begin{remark}[Propriétés avancées]
Les polynômes de Tchebychev forment une base orthogonale de l'espace des fonctions
continues sur $[-1, 1]$ et sont optimaux en norme uniforme parmi les polynômes unitaires
de même degré (résultats admis, cf. théorie de l'approximation).
Pour notre usage dans ChebNet, les deux propriétés suivantes suffisent.
\end{remark}

\begin{theorem}[Propriétés fondamentales des polynômes de Tchebychev]
\label{thm:cheby_properties}
\begin{enumerate}
  \item \textbf{Bornes} : $|T_k(x)| \leq 1$ pour tout $x \in [-1, 1]$ et tout $k \geq 0$.
  \item \textbf{Représentation trigonométrique} : $T_k(\cos\theta) = \cos(k\theta)$.
\end{enumerate}
\end{theorem}

\begin{proof}[Preuve de la propriété 2]
Par récurrence sur $k$ : pour $k = 0$, $T_0(\cos\theta) = 1 = \cos(0)$. \\
Pour $k = 1$, $T_1(\cos\theta) = \cos\theta$. \\
Supposons $T_{k-1}(\cos\theta) = \cos((k-1)\theta)$ et $T_{k-2}(\cos\theta) = \cos((k-2)\theta)$. \\
Alors $T_k(\cos\theta) = 2\cos\theta \cos((k-1)\theta) - \cos((k-2)\theta)$. \\
Or, la formule d'addition donne :
$\cos(k\theta) + \cos((k-2)\theta) = 2\cos\theta \cos((k-1)\theta)$.
Donc $T_k(\cos\theta) = \cos(k\theta)$, ce qui conclut la récurrence.
\end{proof}

\begin{proof}[Preuve de la propriété 1]
Pour tout $x \in [-1,1]$, il existe $\theta \in [0, \pi]$ tel que $x = \cos\theta$.
La propriété~2 donne alors $|T_k(x)| = |\cos(k\theta)| \leq 1$.
\end{proof}

\subsection{Approximation du filtre spectral}

\begin{definition}[Filtre ChebNet]
On approche le filtre spectral $g_\theta$ par sa troncature dans la base de Tchebychev :
\[
  \boxed{g_\theta(\Lt) \approx \sum_{k=0}^{K-1} \theta_k\, T_k(\Lt)}
\]
où $K \in \N^*$ est le nombre de termes retenus et
$\theta = (\theta_0, \ldots, \theta_{K-1}) \in \R^K$
est le vecteur des paramètres apprenables (indépendant de $N$).

\begin{remark}[Convention d'indexation]
On adopte ici la convention du code : $K$ désigne le \emph{nombre total de termes}
(termes $T_0, T_1, \ldots, T_{K-1}$), et non l'ordre maximal.
Ainsi $K = 3$ produit les termes $T_0, T_1, T_2$, ce qui correspond à une réceptivité
à 2 sauts. Certains articles (dont Defferrard \emph{et al.}) utilisent $K$ comme
ordre maximal et écrivent la somme jusqu'à $K$ inclus (soit $K+1$ termes).
\end{remark}
\end{definition}

\begin{remark}
L'argument de $T_k$ est maintenant la \emph{matrice} $\Lt$, et non un scalaire.
La récurrence matricielle correspondante est :
\[
  \bar{\x}_0 = \x, \qquad
  \bar{\x}_1 = \Lt \x, \qquad
  \bar{\x}_k = 2\Lt\, \bar{\x}_{k-1} - \bar{\x}_{k-2},
\]
de sorte que $T_k(\Lt)\x = \bar{\x}_k$.
\end{remark}

\subsection{La convolution ChebNet sans diagonalisation}

\begin{theorem}[Convolution ChebNet]
\label{thm:chebnet_conv}
La convolution ChebNet d'un signal $\x \in \R^N$ avec les paramètres $\theta \in \R^K$ est :
\[
  \boxed{y = g_\theta(\Lt)\, \x = \sum_{k=0}^{K-1} \theta_k\, \bar{\x}_k}
\]
où $\bar{\x}_k = T_k(\Lt)\x$ est calculé par récurrence sans jamais diagonaliser $\L$.
\end{theorem}

\begin{proof}[Élimination de la diagonalisation]
La clé est que chaque produit $\Lt\, \bar{\x}_{k-1}$ est un produit
\textbf{matrice creuse $\times$ vecteur} : $\Lt$ a le même motif de non-zéros que $\A$,
c'est-à-dire exactement $2|\calE| + N$ entrées non nulles (arêtes et diagonale).
Ce produit coûte donc $\mathcal{O}(|\calE|)$ opérations.
En répétant cette récurrence $K$ fois, la complexité totale est $\mathcal{O}(K|\calE|)$,
ce qui ne fait intervenir aucune diagonalisation.
\end{proof}

\subsection{Analyse de complexité comparative}

\begin{table}[h!]
\centering
\caption{Comparaison des complexités : approche stricte vs. ChebNet}
\label{tab:complexity}
\begin{tabular}{@{} l c c @{}}
\toprule
\textbf{Méthode} & \textbf{Complexité temporelle} & \textbf{Paramètres apprenables} \\
\midrule
Spectral strict  & $\mathcal{O}(N^3)$            & $N$ (non transférable) \\
ChebNet ($K$)    & $\mathcal{O}(K |\calE|)$       & $K$ (transférable) \\
\bottomrule
\end{tabular}
\end{table}

\begin{remark}[Localité spatiale]
L'ordre $K$ de l'approximation contrôle directement le \emph{rayon de réceptivité} :
un filtre ChebNet d'ordre $K$ agrège uniquement les informations des voisins
à distance au plus $K$ (au sens de la plus courte distance sur le graphe).
Formellement, si $(i,j)$ est à distance $d(i,j) > K$, alors le chemin de $j$ vers $i$
n'est pas atteint par $K$ multiplications matricielles par $\Lt$,
dont le motif de non-zéros reflète exactement les arêtes de $\calG$.
\end{remark}

\subsection{Généralisation à des signaux multivariés}

En pratique, chaque nœud porte un vecteur de caractéristiques $\x_i \in \R^{F_{\text{in}}}$,
et l'on souhaite produire des sorties dans $\R^{F_{\text{out}}}$.
Le signal est alors une matrice $\mat{X} \in \R^{N \times F_{\text{in}}}$.

\begin{definition}[Couche ChebNet multivariée]
\label{def:chebnet_layer}
Une couche ChebNet est paramétrée par un tenseur de poids
$\Theta \in \R^{K \times F_{\text{in}} \times F_{\text{out}}}$
(un plan $F_{\text{in}} \times F_{\text{out}}$ par ordre polynomial), et sa sortie est :
\[
  \mat{Y} = \sum_{k=0}^{K-1} \bar{\mat{X}}_k\, \Theta_k
  \in \R^{N \times F_{\text{out}}}
\]
où $\bar{\mat{X}}_k = T_k(\Lt)\mat{X} \in \R^{N \times F_{\text{in}}}$
est calculé colonne par colonne via la récurrence de Tchebychev,
et $\Theta_k \in \R^{F_{\text{in}} \times F_{\text{out}}}$ est la matrice de projection
pour l'ordre $k$.

Dans l'implémentation (\texttt{src/layers.py}), ce tenseur correspond exactement au
paramètre \texttt{self.weight} de forme \texttt{(K, in\_features, out\_features)}.
\end{definition}

\begin{remark}[Complexité de la couche multivariée]
La complexité de la couche ChebNet multivariée est :
\[
  \mathcal{O}(K |\calE| F_{\text{in}} + K N F_{\text{in}} F_{\text{out}}).
\]
Le premier terme correspond aux $K$ produits matrice creuse $\times$ matrice dense
(propagation sur le graphe via $\Lt$), et le second aux $K$ projections linéaires finales.
Pour Cora avec $K=3$, $F_{\text{in}} = 1433$, $F_{\text{out}} = 64$ :
la couche d'entrée traite $3 \times 5{,}429 \times 1433 \approx 23\text{M}$ opérations creuses.
\end{remark}

% ============================================================
\section*{Conclusion}
\addcontentsline{toc}{section}{Conclusion}
% ============================================================

Ce document a établi la chaîne mathématique complète menant de la définition d'un graphe
à la couche ChebNet opérationnelle :

\begin{enumerate}
  \item Le \textbf{Laplacien normalisé symétrique} $\L = \I - \D^{-1/2}\A\D^{-1/2}$
    est une matrice semidéfinie positive à spectre dans $[0, 2]$
    (Théorème~\ref{thm:spectrum}), dont les vecteurs propres constituent
    une base de Fourier sur graphe.
  \item La \textbf{convolution spectrale stricte} $g_\theta * \x = \U g_\theta(\Lam)\U^\top \x$
    est exacte mais cubiquement coûteuse $\mathcal{O}(N^3)$
    (Proposition~\ref{prop:complexity_strict}), ce qui la rend inutilisable
    à grande échelle.
  \item L'\textbf{approximation de Tchebychev}
    $g_\theta(\Lt) \approx \sum_{k=0}^{K-1} \theta_k T_k(\Lt)$
    exploite la récurrence à trois termes et la creusité de $\Lt$
    pour descendre à $\mathcal{O}(K|\calE|)$ sans aucune diagonalisation
    (Théorème~\ref{thm:chebnet_conv}, Table~\ref{tab:complexity}).
\end{enumerate}

L'implémentation Python associée (fichiers \texttt{data\_loader.py}, \texttt{graph\_math.py},
\texttt{layers.py}, \texttt{train.py}) traduit exactement ces opérations en tenseurs
PyTorch bruts, en tirant parti de \texttt{torch.sparse.mm} pour les produits creux
et de \texttt{torch.linalg.eigh} pour la version stricte éducative.
Le modèle entraîné sur Cora ($N = 2708$, $K = 3$, 200 époques) atteint
\textbf{78.2\% de précision en test}, sans recourir à aucune bibliothèque GNN.

% ============================================================
\begin{thebibliography}{9}

\bibitem{hammond2011}
D. K. Hammond, P. Vandergheynst, R. Gribonval,
\textit{Wavelets on Graphs via Spectral Graph Theory},
Applied and Computational Harmonic Analysis, 30(2), 129--150, 2011.

\bibitem{defferrard2016}
M. Defferrard, X. Bresson, P. Vandergheynst,
\textit{Convolutional Neural Networks on Graphs with Fast Localized Spectral Filtering},
Advances in Neural Information Processing Systems (NeurIPS), 2016.

\bibitem{kipf2017}
T. N. Kipf, M. Welling,
\textit{Semi-Supervised Classification with Graph Convolutional Networks},
International Conference on Learning Representations (ICLR), 2017.

\bibitem{chung1997}
F. R. K. Chung,
\textit{Spectral Graph Theory},
American Mathematical Society, CBMS Regional Conference Series, Vol. 92, 1997.

\end{thebibliography}

\end{document}
